\chapter{Bnet Instantiations}

\label{ch-bnet-instan}


Random variables are underlined
and their values are not.
For example, $\rva=a$ means
the random variable
$\rva$ takes the value $a$.
A diagram
with all its nodes
underlined represents a
Bayesian Network (bnet),
whereas the same diagram
with the letters not underlined
represents a specific
{\bf instantiation} of that bnet.
For example $\rva\rarrow\rvb\rarrow \rvc$
is a bnet for which
$P(a,b,c)$ equals $P(c|b)P(b|a)P(a)$,
whereas  $a\rarrow b\rarrow c$
is identically equal to  $P(c|b)P(b|a)$. Note that, for convenience,
we define $a\rarrow b\rarrow c$ to
exclude the priors of root nodes such as $P(a)$. An instantiation alone makes no promises.
Think of a {\bf bnet promise} as an equation,
whereas an
 instantiation is just a possible term in an equation.

If $\rva$ is a root node,
then $\PP a$ will denote
$P(a)$.
For example,
\beq
\xymatrix{a&\PP b\ar[l]}
=
P(a|b)P(b)=P(a,b)
\eeq
and 
\beq
\xymatrix{\PP a\ar[r]&b}=
P(b|a)P(a)=P(a,b)
\eeq
More generally,
\beq
\bcen
\xymatrix@R=1pc{
&a_1
\\
\PP x \ar[ru]\ar[rd]
&
\\
&a_2
}\ecen
=
P(x)
\bcen
\xymatrix@R=1pc{
&a_1
\\
x \ar[ru]\ar[rd]
&
\\
&a_2
}\ecen
= 
P(x)P(a_1|x)P(a_2|x)
\eeq


If $\rva$ is a root node,
then $\EE a$ will denote
a weighted sum $\sum_a P(a)$.
For example,

\beq\EE a\rarrow b\rarrow c=\sum_a P(c|b)P(b|a)P(a)
\eeq
More generally,

\beq
\bcen
\xymatrix@R=1pc{
&a_1
\\
\EE x \ar[ru]\ar[rd]
&
\\
&a_2
}\ecen
=
\sum_x P(x)
\bcen
\xymatrix@R=1pc{
&a_1
\\
x \ar[ru]\ar[rd]
&
\\
&a_2
}\ecen
= 
\sum_x P(x)P(a_1|x)P(a_2|x)
\eeq

If $\rva$ is not a root node, then
$\sum a$ signifies
a simple unweighted sum $\sum_a$.
For example,

\beq
x\rarrow\sum a \rarrow y=
\sum_a P(y|a)P(a|x)
\eeq
More generally,
\beq
\bcen
\xymatrix@R=1pc{
&
&\ar[dl]
\\
&
\sum x \ar[lu]\ar[ld]
&
\\
&
&\ar[ul]
}\ecen
=
\sum_x
\bcen
\xymatrix@R=1pc{
&
&\ar[dl]
\\
&
x\ar[lu]\ar[ld] 
&
\\
&
&\ar[ul]
}\ecen=
\sum_x P(\cdot\mid \cdot, x)P(\cdot\mid \cdot, x)P(x|pa(x))
\eeq

\section{Bnet promise and ghost arrows}

A bnet promises a particular full probability distribution
function. We will sometimes refer to that full probability distribution function as the {\bf bnet promise}. 
For example, the bnet 

\beq
\xymatrix{
\rva\ar[r]
&\rvb\ar[r]
&\rvc}
\eeq
promises that

\beq
P(a,b,c)=P(c|b)P(b|a)P(a)
\eeq

Sometimes 
we have to deal with more than one promise at a time.
For example,
in Do Calculus,
the full and amputated graphs
make different promises.
If one doesn't use explicit notation, things
can get confusing when dealing 
with more than one promise 
at a time. 
Suppose

\begin{itemize} 

\item The promise $\pi_1$ equals:

\beq
\begin{array}{c}
P(z, y, x)
=
P(z|y)P(y|z)P(z)
\\
\\
\xymatrix{
\rvz
&\rvy\ar[l]
&\rvx\ar[l]}
\end{array}
\eeq

\item the promise $\pi_2$ equals:

\beq
\begin{array}{c}
Q(z, y, x)=Q(z|y, x) Q(y| x)Q(x)
\\
\\ 
\xymatrix{
\rvz
&\rvy\ar[l]
&\rvx\ar[l]
\ar@/_1pc/[ll]
}
\end{array}
\eeq
\end{itemize}

Note that 

\beq
\begin{array}{l}
\sum_y P(z|y)P(y|x)= P(z|x)
\\
\\
\xymatrix{
z
&\sum y\ar[l]
&x\ar[l]}=
\xymatrix{z&x\ar[l]}
\end{array}
\label{eq-markovian-op-only}
\eeq
and

\beq
\begin{array}{l}
\sum_y Q(z|y, x)Q(y|x)= Q(z|x)
\\
\\
\xymatrix{
z
&\sum y\ar[l]
&x\ar[l]
\ar@/_1pc/[ll]
}=
\xymatrix{z&x\ar[l]}
\end{array}
\label{eq-p-preserving-op}
\eeq
An alternative notation is to
condition on the promise:

\beq
P(\cdot)= P(\cdot\mid \pi_1),
\quad Q(\cdot)=P(\cdot\mid \pi_2)
\eeq

Two bnets are said to be equal iff their promises are equal.
If two bnets are the same except one has more
arrows than the other
(i.e., they have
the same structure in the arrows that are in common, so one
bnet is a \qt{subset}
of the other), we will say the
bnet with more arrows makes a {\bf stronger promise}
and the one with fewer arrows makes a {\bf weaker promise}.
Alternatively, we can say that the weaker bnet has the same number of
arrows as the stronger one, but some have a 0 subscript
 ({\bf ghost arrows}) to
indicate they are {\bf missing}, and thus {\bf carry zero information}.
Graphically,
\beq
\begin{array}{c}
\underbrace{P(x|a_1)}_{
\text{special case  of } P(x|a_1, a_2)
}
P(a_1|pa(a_1))P(a_2|pa(a_2))
= P(x|a_1, \underbrace{a_2}_{\text{ghost}})P(a_1|pa(a_1))
P(a_2|pa(a_2))
\\
\bcen
\xymatrix@R=1pc{
&a_2
&\ar[l]
&
\\
x
&&\ar[dl]
\\
&a_1\ar[ul]
&\ar[uul]
&
}\ecen
=
\bcen
\xymatrix@R=1pc{
&a_2
\ar[dl]_0
&\ar[l]
\\
x
&
&\ar[dl]
\\
&a_1\ar[ul]
&\ar[uul]
}\ecen
\end{array}
\eeq

\beq
P(x|a_1, a_2)= P(x|a_1) \quad \forall a_1, a_2
\eeq
Note that, despite their name, there
is nothing supernatural
about ghost arrows, just
a bigger, more redundant TPM. The TPM $P(x|a_1, a_2)$ is constant 
if we vary $a_2$ while keeping
$(x, a_1)$ fixed.


Note that adding (or removing) ghost arrows
does not change the promise. It's
just a way of expressing the same info with 
a more redundant
(or less redundant) TPM.
Graph operations
that do not change the promise 
will be called {\bf promise-preserving (p-preserving)
operations}.

A result proven for a specific promise is valid for
a weaker promise, but not for a stronger promise.
In other words, a result valid for a bnet is still valid
if an arrow is zeroed (ghosted) but
not if an arrow is unzeroed (unghosted).
Ghost arrows can be added and removed without changing the promise,
but ghost arrows cannot be unghosted.

We often write

\beq
\begin{array}{c}
\sum_b P(c|b)P(b|a)P(a)=P(c,a)
\\
\xymatrix {\PP a\ar[r]&\sum b\ar[r]&c}
\quad=\quad
\xymatrix {\PP a\ar[r]&c}
\label{eq-promi-dep}
\end{array}
\eeq
because  either side summed over $a,c$ yields 1.
This is a p-preserving operation. However, note that
this operation is {\bf promise dependent}; 
i.e., it's
only true if the $P(c,a)$
on the right side
presumes a particular 
\qt{Markovian} promise, namely
\beq
P(c,b, a)=P(c|b)P(b|a)P(a)
\eeq
By comparison,
\beq
\begin{array}{c}
\sum_b P(c|b, a)P(b|a)P(a)=P(c,a)
\\
\xymatrix {\PP a\ar[r]
\ar@/_1pc/[rr]&\sum b\ar[r]&c}
\quad=\quad
\xymatrix {\PP a\ar[r]&c}
\label{eq-promi-indep}
\end{array}
\eeq
is a {\bf promise independent} operation.
\section{Redirecting}

A {\bf fully connected (FC)} bnet (i.e., an acyclic DAG with the maximum possible number of arrows) can be {\bf redirected}
 (i.e., one can reverse the
direction of any arrows 
as long as the final bnet 
is acyclic). This does not change the promise. Graph operations
that do not change the promise 
will be called {\bf
promise-preserving (p-preserving)
operations}.
For example, for a 2 node bnet,
 
\beq
\xymatrix{
\rvx
&\rvy\ar[l]
}=
\xymatrix{
\rvx\ar[r] 
& \rvy
}
\eeq
because both bnets promise 
the same $P(x,y)$. For a 3 node bnet\footnote{Double arrows in an equation means the equation is
true with either choice of arrow
direction.},
\beq
\bcen
\xymatrix{
\rvz\ar[d]\ar[dr]
\\
\rvx\ar@<-1ex>[r]
&\rvy\ar[l]
}\ecen
=
\bcen
\xymatrix{
\rvz\ar[rd]
\\
 \rvx\ar[r]\ar[u]
&\rvy\ar@<-1ex>[ul]
}
\ecen
=
\bcen
\xymatrix{
\rvz\ar[d]
&
\\
\rvx\ar@<1ex>[u]
&\rvy\ar[l]\ar[ul]
}
\ecen
\eeq
because all 6 bnets promise  the
same
 $P(x,y,z)$.


Likewise, 
an instantiation 
of a FC bnet can be redirected without changing the promise.
For example, for a 2 node instantiation,
\beq
P(x, y)=\xymatrix{y&\PP x\ar[l]}
=\xymatrix{\PP y\ar[r]&x}
\eeq
and for a 3 node instantiation,
\beq
P(x,y,z)=
\bcen
\xymatrix{
\PP z\ar[d]\ar[dr]
\\
x\ar@<-1ex>[r]
&y\ar[l]
}\ecen
=
\bcen
\xymatrix{
z\ar[rd]
\\
\PP x\ar[r]\ar[u]
&y\ar@<-1ex>[ul]
}
\ecen
=
\bcen
\xymatrix{
z\ar[d]
&
\\
x\ar@<1ex>[u]
&\PP y\ar[l]\ar[ul]
}
\ecen
\eeq

A useful trick is to
add ghost arrows to an instantiation in order
to complete it and achieve
a FC graph,
and, subsequently, to redirect the {\bf completed}
instantiation. Of course,
none of the arrows of the
redirected graph may be
safely assumed to be ghosts (unless you can prove this via the d-separation theorem,
applied
to the instantiation
prior to redirecting.)
For example,

\beq
\bcen
\xymatrix{
\sum a\ar[dr]
\\
x\ar[u]
&y
}\ecen
=
\bcen
\xymatrix{
\sum a\ar[dr]
\\
x\ar[r]|0
\ar[u]
&y
}\ecen
\\
=
\bcen
\xymatrix{
\sum a
\\
x\ar[r]
\ar[u]
&y\ar[lu]
}\ecen
=
\bcen
\xymatrix{
x\ar[r]&y
}\ecen
\eeq


Recall that a  {\bf clique} 
of a DAG $\calg$ is a 
FC subgraph of $\calg$.
Define an {\bf exploding clique (EC)} as one that
has no incoming arrows,
only outgoing ones.
Redirecting can 
also be done  on ECs
(either bnets or instantiations)
without changing the promise.


\section{P-Preserving Operations}


Recall that we've defined
a p-preserving graph operation  
as one that does not change the promise.
In this section, we shall gather together
all
p-preserving operations
that we have encountered so far.



\begin{claim}
In any graph, a summed leaf node equals 1:

\beq
\begin{array}{c}
\sum_a
P(a|x_1, x_2)f(x_1, x_2, 
\text{anything except $a$})=1
\\
\bcen
\xymatrix@R=.3pc{
x_1\ar[dr]
\\
&\sum a
\\
x_2\ar[ur]
}\ecen=1
\end{array}
\label{eq-collider-sum}
\eeq


\end{claim}
\proof Obvious
\qed

Let us define the
following 3  operations
to be the 
{\bf basic p-preserving
operations}
on a bnet instantiation:
\begin{mybluearea}
\begin{enumerate}
\item
Adding or removing a summed
leaf node.
\item Adding or removing ghost arrows
\item
Redirecting
a FC graph (or EC). This includes adding ghost arrows in order to get an FC graph
to begin with. But after redirecting, 
none of the arrows of the
redirected graph may be
safely assumed to be ghosts (unless you can prove this via the d-separation theorem,
applied
to the instantiation
prior to redirecting.)
\end{enumerate}
\end{mybluearea}

Examples of basic p-preserving operations:
\begin{enumerate}
\item
\beq
\begin{array}{c}
\sum_a P(y|a)P(a)=P(y)
\\
\xymatrix{
\EE a\ar[r]
&y}\quad
=\PP y
\end{array}
\label{eq-diff-priors}
\eeq
Furthermore,

\beq
\begin{array}{c}
\sum_a P(a|y)P(y)=P(y)
\\
\xymatrix{
\sum a
&\PP y\ar[l]}
\quad=\PP y
\end{array}
\label{eq-diff-priors-ee}
\eeq


\item
\beq
\begin{array}{c}
\sum_a P(y|a, x)P(a|x)P(x)=P(y|x)P(x)=P(x,y)
\\
\bcen
\xymatrix@R=1pc{
&\sum a\ar[rd]
\\
\PP x\ar[ru]\ar[rr]
&&y}
\ecen
=
\bcen
\xymatrix{
\PP x\ar[r]
&y}\ecen
\end{array}
\label{eq-med-sum}
\eeq
This follows because the bnet on the left side (LHS) is fully
connected so the arrow $\sum a\rarrow y$ can be reversed and then the sum over $a$ yields 1.

The following identity has a similar justification

\beq
\begin{array}{c}
\sum_a P(y|x,a)P(x|a)P(a)=P(x,y)
\\
\bcen
\xymatrix@R=1pc{
&\EE a\ar[rd]\ar[ld]
\\
x\ar[rr]
&&y}\ecen
=
\bcen
\xymatrix{
\PP x\ar[r]
&y}
\ecen
\end{array}
\label{eq-med-sum-ee}
\eeq

\item
\beq
\begin{array}{c}
\sum_a P(y|a, x_1, x_2)P(a|x_1, x_2)P(x_2|x_1)P(x_1)
= P(y,x_1, x_2)
\\
\bcen
\xymatrix{
\PP x_1\ar[rd]\ar[drr]\ar[dd]
\\
&\sum a\ar[r]
&y
\\
x_2\ar[ru]\ar[rru]
}\ecen
=
\bcen
\xymatrix{
\PP x_1\ar[rrd]\ar[dd]
\\
&&y
\\
x_2\ar[urr]
}
\end{array}
\label{eq-univ-bdoor}
\ecen
\eeq
This follows because the
bnet on the LHS is fully connected, 
so one can reverse the direction of the 
arrow $\sum a \rarrow y$. Then the sum over $a$
yields 1 so the $\sum a$ vertex 
with its 3 incoming arrows can be deleted.

The following identity has a similar justification

\beq
\begin{array}{c}
\sum_a P(y|x_2, a, x_1)P(x_2|a, x_1)P(x_1|a)P(a)
= P(y,x_1, x_2)
\\
\bcen
\xymatrix{
x_1\ar[drr]\ar[dd]
\\
&\EE a\ar[r]\ar[ul]\ar[dl]
&y
\\
x_2\ar[rru]
}\ecen
=
\bcen
\xymatrix{
\PP x_1\ar[rrd]\ar[dd]
\\
&&y
\\
x_2\ar[urr]
}
\ecen
\end{array}
\label{eq-univ-bdoor-ee}
\eeq



\end{enumerate}





\begin{claim}
In an EC, a summed hidden node can be replaced
by a summed non-hidden node:

\beq
\bcen \xymatrix{
*++[F-o]{\EE h}
\ar[d]\ar[dr]
\\
x\ar[r]
&y
}\ecen
=
\bcen \xymatrix{
{\EE z}
\ar[d]\ar[dr]
\\
x\ar[r]
&y
}\ecen
\label{eq-summed-root-identity}
\eeq
and

\beq
\bcen \xymatrix{
*++[F-o]{\sum h}
\ar@{<-}[d]\ar[dr]
\\
\PP x\ar[r]
&y
}\ecen
=
\bcen \xymatrix{
{\sum z}
\ar@{<-}[d]\ar[dr]
\\
\PP x\ar[r]
&y
}\ecen
\label{eq-summed-witness-h}
\eeq
\end{claim}
\proof
\beqa
\bcen \xymatrix{
*++[F-o]{\EE h}
\ar@[red][d]\ar[dr]
\\
x\ar[r]
&y
}\ecen
&=&
\bcen
\xymatrix{
*++[F-o]{\EE h}
\ar@[red][d]\ar[dr]
\ar[r]
&\sum z
\\
x\ar[r]\ar[ur]
&y\ar[u]
}
\ecen
\\
&=&
\bcen
\xymatrix{
*++[F-o]{\sum h}
\ar@{<-}[d]
\ar@{<-}[dr]
\ar@{<-}[r]
&\EE z
\\
x\ar[r]
\ar@{<-}@[red][ur]
&y\ar@{<-}[u]
}
\ecen
\\
&=&
\bcen
\xymatrix{
&\EE z
\\
x\ar[r]
\ar@{<-}@[red][ur]
&y\ar@{<-}[u]
}
\ecen
\eeqa
Eq.(\ref{eq-summed-witness-h})
follows
by reversing the red arrows
in the proof.

\qed